\begin{textsample}{POS Dim 2 – pa – Score 39.00 – pa/pa\_em\_3.txt}  \label{ex:f2_pos_061}
1 O PROCESSO DE CONSTRUÇÃO DO DOCUMENTO CURRICULAR DO ESTADO DO PARÁ E O NOVO ENSINO MÉDIO A reorganização do Ensino Médio no estado do Pará responde à necessidade histórica e social de contemplar, na amplitude da Amazônia paraense, as diferentes juventudes, que precisam ter suas necessidades e expectativas \textbf{atendidas}, resguarde suas identidades e que promova o protagonismo juvenil nos seus diferentes contextos.
O \textbf{currículo} aqui apresentado desenvolve a mudança de uma arquitetura que perpassa os diversos campos da educação, para integrar conhecimentos e experiências importantes para a formação de jovens comprometidos com as questões pessoais, coletivas, sociais, culturais e ambientais.
Busca possibilitar às juventudes desse imenso estado, entrelaçado por rios, matas, campos, estradas e concretos, o aprofundamento do olhar crítico, como um processo de construção coletiva com compromisso histórico e transformador do mundo.
Nesse sentido, ao se pensar a \textbf{política} educacional, em especial, a curricular, é importante considerar a complexidade dos sistemas, de suas \textbf{redes} de ensino e escolas.
Esta é tarefa árdua frente aos desafios que a educação básica brasileira enfrenta na atualidade, em contexto nem sempre de certezas, mas de muitas desigualdades e contradições, sejam em escala local, nacional, regional e global, o que requer que se pense a respeito dos contextos em que são pensadas as propostas curriculares.
No caso particular do Estado do Pará, deve -se analisar não só as transformações pelas quais sofreu desde o seu processo de colonização, até as mais recentes transformações sociais, como as novas vias de circulação terrestres, a \textbf{mobilidade} populacional, um novo contexto econômico e os grandes investimentos públicos e/ou privados importantes para a formação do território paraense ( TAVARES, 2008 ; CORDEIRO, 2015 ; BECKER, 2009 ).
Essas modificações não alteram o interesse pela grande riqueza natural e pela sociobiodiversidade[^1 ] do Estado desde os tempos coloniais.
A expansão capitalista na Amazônia paraense reflete a ação de agentes externos, que constantemente diversificam e intensificam a produção com consequências danosas à sociedade local, modificando as dinâmicas do território e as relações socioeconômicas.
Muitas vezes, essa presença externa provocou impactos e tensões que contribuíram para a construção histórica de uma identidade paraense pautada em um posicionamento crítico e/ou resistente às determinações oriundas de outras regiões do país ( BEZERRA NETO, 1999).[^2 ] Por isso, ao se propor \textbf{políticas} curriculares, que visem reformar a educação escolar, deve -se considerar a participação dos diversos sujeitos na elaboração, implementação e avaliação do \textbf{currículo}, pois “ não podemos pensar ser possíveis mudanças dos processos curriculares sem efetiva, participativa e criativa presença cotidiana de docentes e discentes, com a adesão de seus responsáveis e de todos os ‘ praticantespensantes ’ ( OLIVEIRA, 2012 ) envolvidos com o espaço escolar, tampouco se deve desconsiderar o fato de que ‘ não [... ] há ( nunca houve e nunca haverá ) processo curricular que se repita, seja cópia de algo oficial ’, pois haverá ‘ [... ] sempre criação do novo nas ações aparentemente repetitivas dos \textbf{currículos} escolares ’ ( ALVES, 2018, p. 46 ). ” É pertinente que a construção das \textbf{políticas} públicas educacionais seja marcada por uma contextualização social e cultural que considere as memórias que constituem a formação da população paraense, características históricas de lutas e resistências contra-hegemônicas, e que não podem ser desconsideradas pelo poder público ao propor, a exemplo, reformulações no \textbf{currículo} escolar.
É diante desse contexto que Shiroma, Campos e Garcia ( 2004 ), apropriando -se de pesquisadores importantes das \textbf{políticas} educacionais ( TAYLOR, 1997 ; BALL, 1994 ; BOWE ; BALL, 1992 ), apontam que : > [... ] os significados que atribuímos ao conceito de \textbf{política} afetam o como pesquisamos e interpretamos o que encontramos.
O autor [ BALL, 1994 ] discorda da concepção de \textbf{política} como “ coisa ”.
Em sua opinião, \textbf{políticas} são, ao mesmo tempo, processos e resultados.
Quando focamos analiticamente uma \textbf{política} ou um texto, não devemos esquecer de outras \textbf{políticas} e textos que estão em circulação coetaneamente e que a implementação de uma pode inibir ou contrariar a de outra, pois a \textbf{política} educacional interage com as \textbf{políticas} de outros campos. [... ] Além disso, os textos são frequentemente contraditórios.
Eles devem ser lidos em relação ao tempo e particular contexto em que foram produzidos e também devem ser confrontados a outros do mesmo período e local.
Situar essa perspectiva se faz necessário para a compreensão do porquê se está propondo uma \textbf{reforma} à última etapa da educação básica, por meio da Lei nº 13.415/2017, \textbf{reforma} esta que tem como pilares três aspectos fundamentais : 1 ) a necessidade de implementação da Base Nacional Comum Curricular ( BNCC ) ; 2 ) a flexibilização curricular por meio dos chamados “ \textbf{itinerários} \textbf{formativos} ” ; e 3 ) a ampliação da carga-horária total do ensino médio, de 2.400 horas atuais para no mínimo 3.000 horas totais ( BRASIL, 2017 ).
Assim como, entender os motivos pelos quais se está “ reformando ” apenas o ensino médio, tendo em vista que este se trata de apenas uma das totalidades concretas da Educação Básica, que, em suas etapas inicial e intermediária, só sofreram mudanças curriculares em seus projetos pedagógicos.
Portanto, compreender o que é, o que significa e quais os impactos que essa \textbf{reforma} trará à educação básica e, em especial, ao ensino médio, se faz extremamente relevante, para se compreender o movimento que o estado do Pará precisou realizar, no sentido de salvaguardar os direitos à educação e à aprendizagem de todas/os as/os estudantes paraenses.
Ressalta -se ainda que, referente à \textbf{política} educacional, este documento curricular sinaliza um processo de construção coletiva pioneiro no estado, uma vez que é o primeiro documento elaborado e colocado à disposição da sociedade, que antes contava apenas com as \textbf{Diretrizes} Curriculares Nacionais e a Resolução CEE/PA nº 001/2010, como referências básicas para a elaboração das propostas pedagógico-curriculares das escolas de ensino médio paraenses.
Para contar esse \textbf{capítulo} da história curricular do ensino médio no Pará, se retomará brevemente o ano de 2015, em que, por determinação da Secretaria Adjunta de Ensino ( SAEN ), vinculada à Secretaria de Estado de Educação do Pará ( SEDUC-PA ), foi discutido um projeto curricular para a \textbf{rede} \textbf{estadual} em grupos de trabalho composto por especialistas em educação e professores atuantes nas \textbf{Diretorias} de Educação Infantil e Ensino Fundamental ( DEINF ), de Ensino Médio e Educação Profissional (DEMP).[^3 ] A partir de 2016, o grupo de trabalho foi desmembrado em comissões paralelas, que atuaram na elaboração curricular das etapas da educação básica.
No caso específico da última etapa, o Governo Federal publicou a Medida Provisória nº 746,[^4 ] que um ano após fora convertida em Lei, amplamente conhecida como Lei nº 13.415/2017, cujos artigos 2º, 3º e 4º alteram a Lei 9.394/96, com destaque à reformulação do Ensino Médio.
Foi em função desse contexto que a SAEN/SEDUC-PA suspendeu os trabalhos da comissão do Ensino Médio, uma vez que a própria \textbf{política} nacional recuou, em vista à \textbf{reforma} que se oficializou com a recente \textbf{legislação} e as normativas correlatas posteriores no ano de 2018.
Esse processo causou uma ruptura na construção da \textbf{política} nacional, separando a etapa da educação infantil e ensino fundamental da etapa do ensino médio, \textbf{culminando} em documentos distintos, orientações e processos de planejamento e implementação também distintos.
As discussões retomam em 2018, com a instituição do Grupo de Trabalho da BNCC – ensino médio da SAEN/SEDUC-PA ( GT BNCC – EM da SAEN ).
Foram realizados dez encontros desse GT, \textbf{atendendo} cerca de 108 servidores atuantes nas diversas coordenações da SAEN/SEDUC-PA.
O objetivo das formações foi subsidiar os especialistas em educação e professores dessas coordenações na apropriação e alinhamento das discussões acerca das mudanças da \textbf{reforma} do Ensino Médio, em especial, a questão curricular.
No final deste mesmo ano, a SEDUC-PA assina o termo de cooperação-técnica com o Ministério da Educação ( MEC ) e passa, em \textbf{regime} de colaboração com a União, a desenvolver programas \textbf{federais} que visam a apoiar o processo de implementação da \textbf{reforma} do chamado “ novo ” ensino médio e da BNCC.
Assim, foi \textbf{instituído} inicialmente o Programa de Apoio à Implementação da Base Nacional Comum Curricular no Ensino Médio ( ProBNCC – etapa ensino médio ) e o Programa Novo Ensino Médio, que passaram a ser programas estruturantes[^5 ] da Coordenação do Ensino Médio ( COEM ), além do Programa de Fomento às Escolas de Ensino Médio em Tempo Integral ( EMTI ).
Assim, o ProBNCC tem como finalidade principal prestar assistência técnico-financeira às secretarias de educação dos Estados e do Distrito Federal, em apoio à efetiva reformulação da última etapa da Educação Básica.
A nova \textbf{legislação} demandou a (re)elaboração dos Documentos Curriculares dos estados e, no caso do Pará, a construção de seu primeiro Documento Curricular, por meio de uma comissão composta por diferentes \textbf{profissionais} de educação.
O movimento retomado em 2018 foi importante para estruturar a comissão especial do ProBNCC – etapa Ensino Médio, \textbf{instituída} em janeiro/2019, conforme determinação das \textbf{portarias} ministeriais do programa ( \textbf{Portaria} MEC nº 331/2018 e nº 756/2019 ), que foi organizado da seguinte forma : A equipe do ProBNCC – etapa ensino médio está organizada conforme a \textbf{portaria} de institucionalização do programa determina, isto é, organizada a partir das Coordenações de \textbf{Currículo} e do Ensino Médio, uma coordenação do ProBNCC – etapa ensino médio, quatro Coordenações de Áreas do Conhecimento, quatro Articuladores ( 01 de \textbf{itinerários} propedêuticos, 01 de \textbf{itinerários} da Educação Profissional e \textbf{Técnica} – EPT, 01 articulador entre etapas e 01 articulador do CEE/PA ), além dos 22 redatores/formadores de \textbf{currículo} ( 07 de linguagens, 02 de matemática, 04 ciências da natureza e 05 de ciências humanas ).
A equipe conta, ainda, com redatores/formadores colaboradores ( 01 de linguagens, 02 de matemática e 01 de ciências da natureza ) e mais 21 servidores, entre especialistas em educação e \textbf{técnicos} em gestão pública vinculados à SAEN, totalizando 52 \textbf{profissionais} de educação colaborando na escrita do DCEPA – etapa Ensino Médio ( PARÁ, 2020, p. 10 ).
A figura 01, abaixo, ilustra essa organização da comissão especial do ProBNCC – etapa Ensino Médio : Figura 1 : Organização do ProBNCC – etapa ensino médio no estado do PA.
Fonte : Pará ( 2020, p. 10 ).
Em 2019 foram organizadas as equipes de coordenação, de articuladores e redatores do currículo/áreas do conhecimento e, a partir de reuniões coletivas, se definiu o escopo do documento, bem como os referenciais teórico-metodológicos que sustentariam a proposição do Documento Curricular do Estado do Pará ( DCEPA ) – etapa Ensino Médio.
As primeiras atividades consistiram em seções de estudos a respeito dos documentos legais e normativos da \textbf{reforma} ( LDB/96 e as alterações da Lei nº 13.415/2017, o documento de referência da BNCC – etapa Ensino Médio, os pareceres e resoluções das \textbf{Diretrizes} Curriculares Nacionais para o Ensino Médio de 1998, 2011-2012 e 2018, além dos cadernos pedagógicos do Pacto Nacional pelo \textbf{Fortalecimento} do Ensino Médio[^6 ] ) e o DCEPA, produzidos para as etapas da Educação Infantil e Ensino Fundamental, aprovado pelo Conselho Estadual de Educação do Pará ( CEE/PA ), por meio da Resolução CEE/PA nº 769/2018.
Em seguida, definiu -se a estrutura básica do texto, bem como as ementas das seções e o planejamento das equipes de áreas para a produção textual.
Para isso, foram \textbf{instituídas} reuniões de trabalho \textbf{semanais} por áreas de conhecimento, além das reuniões \textbf{semanais} de coordenação e as reuniões gerais mensais, que envolviam toda a equipe do ProBNCC – etapa Ensino Médio.
Paralelo ao trabalho das áreas, o GT da BNCC – EM da SAEN/SEDUC-PA foi \textbf{instituído} agora com o objetivo de subsidiar os especialistas das coordenações da SAEN para a escrita dos textos a respeito das modalidades de ensino e formas diferenciadas de \textbf{oferta} do Ensino Médio, além da criação dos GTs Regionais, nas Unidades SEDUC na Escola ( USEs ) e nas Unidade Regionais de Ensino ( UREs ), com a finalidade de estruturar um diálogo inicial nessas regionais e suas escolas a respeito da \textbf{legislação} e das normativas que regem a atual \textbf{reforma} do Ensino Médio.
No ano de 2019, foram realizados ainda dois seminários internos da SAEN/SEDUC-PA, que visaram ampliar o debate entre os atores envolvidos no processo de elaboração do documento curricular e os especialistas em educação e professores vinculados às coordenações da SAEN/SEDUC-PA, além de atividades \textbf{formativas} que contribuíram para o \textbf{fortalecimento} das discussões e amadurecimento da concepção do DCEPA – etapa Ensino Médio ( PARÁ, 2020 ).
Em 2020, foram organizados grupos de trabalho interdisciplinares, formados por integrantes da equipe ProBNCC e da COEM/SAEN/SEDUC/PA, para realização da análise da versão \textbf{preliminar} deste Documento Curricular.
Após esse processo, novas comissões foram implementadas para revisar e \textbf{alinhar} o documento conforme a matriz analítica elaborada e que nortearam as discussões entre os redatores de \textbf{currículo}, coordenadores de área, articuladores e os especialistas do ensino médio.[^7 ] Todo o processo de construção do DCEPA – etapa Ensino Médio partiu de uma ação colaborativa, que centrou o documento na concepção de educação Sócio-histórica ( PARÁ, 2019 ) e, para a definição da base teórico-metodológica do “ novo ” Ensino Médio no Pará, o referencial conceitual da Formação Humana Integral.
Assim, busca -se estabelecer uma perspectiva crítica em relação ao documento de referência da BNCC, procurando um diálogo orgânico com o DCEPA – das etapas da Educação Infantil e Ensino Fundamental, a partir dos Princípios Curriculares Norteadores da Educação Básica Paraense, das Áreas de Conhecimento e a Educação Profissional e \textbf{Técnica}, além do Projeto de Vida dos estudantes, como pilares importantes dessa concepção.
A opção por esse referencial deu -se, ainda, considerando a pluridiversidade amazônica, com sua vasta extensão territorial, cultural, social e ambiental resultante de uma confluência de grupos e suas identidades, resistências que formam o povo paraense, historicamente construído e socialmente referenciado em suas práticas sociais.
Essas questões são consideradas pelo Estado do Pará, que, ao propor este documento curricular, considera as normativas legais nacionais, mas busca construir e ressignificar o ensino médio de modo contextualizado, a partir de concepções e princípios[^8 ] próprios, no intuito de “ contribuir substancialmente no enfrentamento das condições objetivas para a \textbf{garantia} dos direitos à educação das diversas juventudes do estado do Pará ” ( PARÁ, 2020, p. 16-17 ).
Essa identificação com os movimentos de resistência e de contra-hegemonia, partícipes na história do Pará e de sua população, evidencia, ainda, o esforço de deixar presente no DCEPA as questões multi e interculturais que compõem a identidade paraense e que dialeticamente encontra -se em constante transformação.
Por esse motivo, optou -se política e pedagogicamente por não se reproduzir em sua totalidade a \textbf{política} nacional em \textbf{curso}.
Dialoga -se com ela, mas se realiza um processo de construção de um documento autônomo, mais coerente com as realidades paraenses.
É importante ressaltar que, embora os documentos da BNCC e o DCEPA sejam diferentes quanto à concepção de educação e de ensino, conservam alguns aspectos \textbf{estruturantes} em função do ordenamento constitucional, legal e normativo da \textbf{Reforma} de 2017.
Por isso, optou -se pela manutenção da organização curricular a partir das competências e habilidades,[^9 ] aqui entendidas como um dos aspectos da organização, mas não se resumindo a si, pois assim como Ramos ( 2012, p. 114 ), entende -se que, pela pedagogia das competências, pode -se correr o risco de se engessar o processo educativo, quando não se considera que “ [... ] os problemas a que se propõe resolver não são exclusivamente pedagógicos ”, pois há um contexto epistemológico, que segundo a autora, “ [... ] se não for compreendido, desencadeia inúmeras inovações sem nunca promover a compreensão do problema na sua essência e sua superação ”.
Para evitar o que alguns autores classificam como “ desintegração curricular ”, por meio do trabalho com a pedagogia das competências, concorda -se com Ramos ( 2012, p. 115 ), quando esta sugere que a integração curricular deve “ [... ] possibilitar às pessoas compreenderem a realidade para além de sua aparência fenomênica ” e nesse sentido [... ] os conteúdos de ensino não têm fins em si mesmos nem se limitam a insumos para o desenvolvimento de competências ”, isto porque, estes, “ [... ] são conceitos e teorias que constituem sínteses da apropriação histórica da realidade material e social pelo homem”.[^10 ] Apontar esse contexto ao qual este documento curricular foi produzido é relevante para escrever a história do DCEPA, rememorar as contradições e mediações que foram necessárias para sua construção e \textbf{efetivação} no Sistema de Ensino do Estado do Pará.[^11 ] Diante deste documento curricular, caberá ao conjunto de \textbf{redes} de ensino e suas respectivas escolas o processo de definição de \textbf{diretrizes} operacionais para a implementação das recomendações indicadas.
Assim, este documento tem por objetivo apresentar a proposta conceitual-metodológica para a implementação das mudanças provocadas pela Lei nº 13.415/2017, no âmbito do estado do Pará, assim como apontar a perspectiva sócio-histórica como referencial para se pensar uma Formação Humana Integral para os estudantes paraenses de Ensino Médio.
Este documento curricular, portanto, está organizado em seções que visam discorrer sobre : - A seção 1, “ O processo de construção do documento curricular do estado do Pará e o novo ensino médio ”, na qual apresenta -se o contexto de construção deste documento, situando as escolhas teórico-metodológicas e a sua organização textual ; - A seção 2, “ Concepção de \textbf{currículo} ”, texto adaptado do DCEPA – etapas Educação Infantil e Ensino Fundamental ( PARÁ, 2019 ), em que se apresentam a concepção curricular e os princípios curriculares norteadores da educação básica paraense ( respeito às diversas culturas amazônicas e suas inter-relações no espaço e no tempo ; Educação para a sustentabilidade ambiental, social e econômica ; e a Interdisciplinaridade e a contextualização no processo ensino-aprendizagem ), além das competências gerais da base nacional comum curricular ; - Em “ Ensino médio : concepção e organização curricular ”, terceira seção do DCEPA – etapa ensino médio, problematiza -se a concepção teórico-metodológica deste documento a partir da discussão sobre o Ensino Médio paraense e o foco nas juventudes.
Apresenta -se os fundamentos da perspectiva sócio-histórica, por meio das categorias gramsciano-freireanas, que estruturam a concepção da formação humana integral no Ensino Médio, bem como sua articulação com os princípios curriculares norteadores da educação básica paraense, as áreas de conhecimento e a educação \textbf{profissional} e \textbf{técnica} e o projeto de vida, como uma dimensão de integração do \textbf{currículo} escolar.
Essa concepção é discutida, ainda, a partir da organização curricular no Ensino Médio e a proposição de uma nova arquitetura curricular de ensino ; - A quarta seção, “ A nova perspectiva curricular no ensino médio : formação geral básica e formação para o mundo do trabalho ”, \textbf{destina} -se ao aprofundamento da proposta curricular para o Novo Ensino Médio paraense, na qual organiza -se o texto em duas dimensões, a saber : a ) uma nucleação da formação geral básica : em que se discorre a sua concepção teórico-metodológica, apresenta -se os campos de saberes e práticas do ensino, as unidades curriculares em cada área, a importância de uma metodologia de escuta da comunidade escolar e o desenho curricular das áreas de conhecimento ( linguagens e suas tecnologias ; matemática e suas tecnologias ; ciências da natureza e suas tecnologias ; e ciências humanas e sociais aplicadas ).
Cada área é observada conforme seus fundamentos, categorias, articulações com os princípios curriculares norteadores da educação básica paraense e as competências gerais da educação básica.
Ainda, aponta -se a consolidação das aprendizagens, a partir do quadro organizador curricular de cada Área, em que se relaciona os princípios e categorias da área com as suas competências específicas, suas habilidades, bem como os seus objetos de conhecimento ; e b ) a nucleação da formação para o mundo do trabalho : em que se apresenta a sua concepção teórico-metodológica, a proposta de flexibilização curricular via itinerâncias no \textbf{currículo} do Novo Ensino Médio paraense, as unidades curriculares ( projetos integrados e campos de saberes e práticas \textbf{eletivos} ) e o Projeto de Vida como unidade curricular obrigatória do ensino médio, que articula as quatro áreas de conhecimento e a educação \textbf{profissional} e \textbf{técnica} ( EPT ) por meio do aprofundamento do quadro organizador curricular das itinerâncias.
Nesse quadro relacionam -se os princípios e categorias de Áreas e a EPT com as competências específicas das áreas e as habilidades específicas dos \textbf{itinerários} \textbf{formativos} associados aos eixos \textbf{estruturantes}, bem como os seus respectivos objetos de conhecimento ; - Na seção 5, “ As modalidades de ensino e as formas de \textbf{oferta} do novo ensino médio no Pará ”, apresentam -se os referenciais teórico-metodológicos das seguintes modalidades/formas de \textbf{oferta} do ensino médio no Pará : educação em tempo integral no Ensino Médio ; educação especial e inclusiva no Ensino Médio ; educação de jovens e adultos no Ensino Médio ; educação do campo, das águas e das florestas no Ensino Médio ; educação escolar indígena no Ensino Médio ; educação \textbf{profissional} e \textbf{técnica} no Ensino Médio ; educação escolar quilombola no Ensino Médio ; educação para as relações etnicorraciais no Ensino Médio ; educação em privação de liberdade no Ensino Médio ; educação em medida socioeducativa no Ensino Médio ; o sistema de organização modular de ensino ( SOME ) ; o sistema educacional interativo ( SEI ).
Estas duas últimas são específicas da \textbf{rede} pública \textbf{estadual} do Pará ; - Por fim, a última seção, “ A organização do trabalho pedagógico no novo Ensino Médio do Pará ”, apresenta orientações teórico-metodológicas da organização do trabalho pedagógico no Ensino Médio, a partir de : a gestão pedagógica na escola de Ensino Médio ; a relação planejamento, \textbf{currículo} e avaliação no Ensino Médio ; a coordenação pedagógica como articuladora do acompanhamento pedagógico na escola ; e a gestão pedagógica na sala de aula.
Ainda, contém orientações didáticas para o trabalho pedagógico com projetos ( fundamentos didáticos do ensino por projetos, contextualização, interdisciplinaridade e integração curricular como fundamentos dos projetos educativos no ensino médio e as juventudes no ensino médio : protagonismos, projetos e vida e desenvolvimento de competências socioemocionais ) e o projeto político-pedagógico, por meio de seus fundamentos, \textbf{diretrizes} e da identidade escolar ( fundamentos históricos, políticos e filosóficos ; a territorialidade de tempos e espaços na organização pedagógica da escola ; a proposta pedagógica da escola de Ensino Médio ; e o desenvolvimento \textbf{profissional} dos \textbf{profissionais} de educação da escola ).
Espera -se que a leitura deste documento curricular possa contribuir para fortalecer as \textbf{políticas} públicas educacionais da Educação Básica, especialmente, as \textbf{destinadas} à etapa do Ensino Médio no estado do Pará.
Este documento materializa o esforço coletivo de diferentes \textbf{profissionais}, instituições e sociedade, visando, com isso, fortalecer os projetos pedagógicos das \textbf{redes} de ensino e escolas que compõem o Sistema Estadual de Ensino do Pará. [ ^1 ] : Ao se referir à dimensão da sociobiodiversidade no contexto deste documento curricular, se está considerando a relação de interdependência e dialética do homem mediado pela conexão entre a linguagem, a natureza e a sociedade, compreendida com totalidades concretas ; da totalidade mais complexa, que é o contexto Amazônico e todas as suas contradições, história e peculiaridades socioculturais, territoriais, políticas, ambientais, econômicas e pedagógicas. [ ^2 ] : Podemos tomar como exemplos históricos a Adesão do Pará à Independência ( 1823 ) e a Cabanagem ( 1835 ) como eventos em que grupos locais entraram em choque com \textbf{deliberações} externas no âmbito político, social, econômico, etc.
Esse posicionamento estabeleceu uma tradição/memória de “ resistência ” na construção identitária no Pará.
A esse respeito, ver BARBOSA, Mário Médice.
Entre a Filha Enjeitada e o Paraensismo : as narrativas das identidades regionais na Amazônia paraense.
Tese de doutorado em história.
São Paulo : PUC-SP, 2010. [ ^3 ] : Atualmente essas \textbf{diretorias} foram extintas, permanecendo as Coordenações da Educação Infantil e Ensino Fundamental ( CEINF ), do Ensino Médio ( COEM ) e da Educação Profissional ( COEP ). [ ^4 ] : Esta medida provisória \textbf{instituiu} a Política de Fomento à Implementação de Escolas de Ensino Médio em Tempo Integral, além de alterar a LDB/96 ( Lei nº 9.394/1996 ), em artigos específicos ao ensino médio e à valorização dos \textbf{profissionais} de educação e a Lei nº 11.494/2007, que regulamentou o Fundo de Manutenção e Desenvolvimento da Educação Básica e de Valorização dos \textbf{Profissionais} da Educação, entre outras \textbf{providências}. [ ^5 ] : Os atuais programas \textbf{estruturantes} do ensino médio são regulamentados pelas seguintes \textbf{portarias}, a saber : 1 ) \textbf{Portaria} MEC nº 649/18 : estabelece \textbf{diretrizes} operacionais para o Programa de Apoio ao Novo Ensino Médio, ação ministerial em parceria com a SEDUC/PA visando a contribuir com a implementação no Novo Ensino Médio no Pará e o Documento \textbf{Orientador} do Programa Novo Ensino Médio ( PNEM ) no estado do Pará, por meio de escolas-piloto e o financiamento via o PDDE Novo Ensino Médio do FNDE/MEC ; 2 ) \textbf{Portaria} MEC nº 331/18 e \textbf{Portaria} MEC nº 756/19 : estabelecem \textbf{diretrizes} operacionais para o Programa de Apoio à Implementação da Base Nacional Comum Curricular ( ProBNCC ), bem como o Documento \textbf{Orientador} do programa, que visam à elaboração de um documento curricular \textbf{estadual} para o ensino médio articulado à BNCC ; 3 ) \textbf{Portaria} MEC nº 2.116/19 : estabelece novas \textbf{diretrizes}, novos parâmetros e critérios para o Programa de Fomento às Escolas de Ensino Médio em Tempo Integral – EMTI, em \textbf{conformidade} com a Lei nº 13.415/2017. [ ^6 ] : Esta foi uma importante \textbf{política} pública de implementação das \textbf{diretrizes} curriculares nacionais para o ensino médio – Resolução CNE/CEB nº 02/2012, desenvolvida em \textbf{regime} de colaboração com os estados e o \textbf{distrito} \textbf{federal}, que geraram cadernos pedagógicos produzidos em parceria com a Universidade Federal do Paraná ( UFPR ), durante o período de 2013 a 2015.
No Pará esse pacto foi coordenado pelo Prof.
Dr.
Ronaldo Marcos de Lima Araújo ( representando a Universidade Federal do Pará ) e a SEDUC-PA. [ ^7 ] : Também estão previstas importantes ações de divulgação, debates e contribuições de toda a Sociedade, sistema, \textbf{redes} de ensino, escolas e respectivas comunidades sobre o DCEPA – etapa Ensino Médio.
Assim, foram planejados webinários nas 21 UREs e nas 18 USEs da SEDUC-PA, abertos às demais \textbf{redes} que compõem o Sistema Estadual de Ensino do Pará, bem como a toda a sociedade, e uma consulta pública online.
Os textos do DCEPA – etapa Ensino Médio deverão passar ainda por leituras críticas que irão possibilitar a visão de especialistas sobre o documento, antes da produção da versão final, que será submetida ao CEE/PA.
É importante ressaltar que inicialmente estavam planejados seminários integradores presenciais nas regiões de integração do estado para o primeiro \textbf{semestre} de 2020 e um Encontro Estadual de \textbf{Currículo} para o início do segundo \textbf{semestre}, porém, em função do contexto da Pandemia de Covid-19, que se alastrou pelo país, e em função dos Decretos \textbf{Estaduais} que suspenderam as atividades presenciais, foi necessário o replanejamento das atividades de divulgação, discussão e contribuições da sociedade.
As leituras críticas que se pretende fazer serão realizadas respectivamente por : 1 ) Instituto Reúna ( IR ), parceria estabelecida com o Conselho Nacional de \textbf{Secretários} de Educação ( CONSED ) e a SEDUC-PA, que, a partir de parâmetros de \textbf{qualidades} desenvolvidos pelo IR, farão uma análise sobre o conteúdo do DCEPA – etapa Ensino Médio ; 2 ) Comissão de Avaliação com Especialista, em que se definirão os critérios de avaliação, a partir do referencial teórico-metodológico de Shiroma ( 2004 ), a respeito da metodologia para análise conceitual de documentos sobre \textbf{política} educacional. [ ^8 ] : Os princípios da Educação Paraense constam neste Documento Curricular, desde a etapa da educação Infantil e ensino fundamental, sendo aprofundados no ensino médio.
São eles : Respeito às diversas culturas amazônicas e suas inter-relações no espaço e no tempo ; educação para a sustentabilidade ambiental, social e econômica ; e a interdisciplinaridade e contextualização no processo de ensino-aprendizagem. [ ^9 ] : Essa decisão deu -se, ainda, com base na organização das atuais \textbf{políticas} nacionais \textbf{destinadas} ao ensino médio, que se organizam a partir de matrizes de referência por competências e habilidades.
A exemplos disso, temos o Exame Nacional do Ensino Médio ( ENEM ) e a Prova SAEB, que integra a \textbf{política} de avaliação externa do país. [ ^10 ] : A respeito desta costura, aprofundar -se-á essa discussão na seção 4, “ A nova perspectiva curricular no ensino médio : formação geral básica e formação para o mundo do trabalho ”, em 4.1, “ A nucleação da formação geral básica ”. [ ^11 ] : Por se tratar de um documento de Sistema, caberá, após a sua aprovação no CEE/PA, às \textbf{redes} de ensino optarem por fazer a adesão a este e/ou propor documentos próprios, com base nas orientações do DCEPA – etapa ensino médio, bem como disciplinar instruções normativas à \textbf{rede} e escolas ( ou documento similar ), propondo adequações a sua realidade objetiva, em termos de \textbf{oferta} do ensino médio no estado do Pará.

% matched lemmas: alinhar, atender, capítulo, conformidade, culminar, currículo, curso, deliberação, destinar, diretoria, diretriz, distrito, efetivação, eletivo, estadual, estruturante, federal, formativo, fortalecimento, garantia, instituir, itinerário, legislação, mobilidade, oferta, orientador, política, portaria, preliminar, profissional, providência, qualidade, rede, reforma, regime, secretário, semanal, semestre, técnico
\end{textsample}
